\documentclass[12pt]{article}
\usepackage[utf8]{inputenc}
\usepackage[T1]{fontenc}
\usepackage{amsmath,amssymb}
\usepackage[backend=biber]{biblatex}
\addbibresource{refs.bib}

\title{Parallel Computing Proposal: MAKAsolve}
\author{
	Mikiel Gica
	\and Aidan Hoover
	\and Kiran Koch-Mathews
	\and Aiden Woodruff
}
\date{March 27, 2026}

\begin{document}
\maketitle

Computational Fluid Dynamics (CFD) involves solving for flow variables using
the Navier-Stokes (NS) equations on a domain, typically discretized by a
mesh. Solving the NS equations using numerical methods such as Finite Element
or Finite Volume can become computationally expensive and intractable to
solve on a single processor.  This expensive operation can be parallelized by
splitting (partitioning) the mesh into parts that one processor can handle.
Each processor can then solve the NS equations on its own partition.

The NS equations can be written as a system of multiple advection-diffusion
equations. To emulate the process of solving the NS equations on a
partitioned mesh, we want to solve a single advection-diffusion equation
using Finite Element methods. One form of the advection-diffusion equation
for a conserved variable \(\phi\) is

\begin{equation}
	\begin{split}
		\phi_{,t} + a_i\phi_i - \kappa\phi_{,ii} = f & \text{ on } \Omega \\
		\phi = g & \text{ on } \Gamma_g \label{eq:strong1-ess} \\
		-\kappa\phi_{,i} = h & \text{ on } \Gamma_h
	\end{split}
\end{equation}

where \(a_i\) are the components of an advection vector, \(\kappa > 0\) is a
diffusion constant, \(f : \Omega \to \mathbb{R}\) is a source term, \(g\) is
the essential (Dirichlet) boundary condition, \(h\) is the natural (Neumann)
boundary condition, and \(i\) is a variable from 1 to the number of spatial
dimensions (Einstein notation is used for all equations). We make the
following simplifying assumptions: \(\phi_{,t} = 0\) (steady state),
\(f = 0\) (no source term), and \(h = 0\) (only Dirichlet BCs, i.e.
\(\Gamma_h = \empty\). The Dirichlet BC is satisfied a priori by selecting a
function
\(\phi \in \delta = \left\{\phi | \phi \in H^1, u = g |_\Gamma\right\}\). The
space \(H^1\) is the set of functions which are continuous and
square-integrable
\(\int_\Omega\left(\phi^2 + \phi^2_{,x}\right)d\phi < \infty\) up to the 1st
derivative (so that the solution does not blow up). Then the simplified
strong form PDE is

\begin{equation}
	\begin{split}
		a_i\phi_i - \kappa\phi_{,ii} = 0 & \text{ on } \Omega \\
		\phi = g & \text{ on } \Gamma
	\end{split}
\end{equation}

To ease calculation of the solution, the method of weighted residuals and
integration by parts are applied to obtain the weak form that can be solved
instead:

\begin{equation}
	\int_\Omega w_{,i}\left(-a_i\phi + \kappa\phi_{,i}\right) d\Omega = 0
\end{equation}

The weighting function \(w : \Omega \to \mathbb{R}\) is also in the set
\(H^1\) to ensure that the solution does not blow up. To solve the solution
in the mesh, the weak form is converted to the Galerkin form
\cite{hughes-1987}:

\begin{equation}
	\int_\Omega \bar{w}_{,i}(-a_i\bar{\phi} + \kappa\bar{\phi}_{,i}) d\Omega
	= 0
\end{equation}

Then the Galerkin form can be used to form a stiffness matrix which can be
solved with typical linear algebra solvers. Since both \(phi\) and \(w\) are
\(H^1\), linear finite elements are used later to discretize the solution
funciton. However, the Galerkin formulation is known to be unstable in
advection dominated cases \cite{hughes-2010}. One method, known as Stabilized
Upwind Petrov Galerkin (SUPG), introduces a stabilization term proportional
to the strong residual. The resulting weak form residual is given below:

\begin{equation}
	\begin{split}
		\int_\Omega
			\bar{w}_{,i}(-a_i\bar{\phi} + \kappa\bar{\phi}_{,i}) d\Omega
		+ \sum_{e=1}^{N_e}
			\int_{\Omega^e} (-a_i\bar{\phi}_{,i})(-\tau R(\bar{\phi})) d\Omega
		= 0 \\
		R(\bar{\phi}) = a_i\bar{\phi}_i-\kappa\bar{\phi}_{,ii}
	\end{split}
\end{equation}

The left factor of the stabilization term has no diffusive component because
linear finite element are used, and since we are careful to only integrate
over element boundaries, the \(\phi_{,ii} = 0\). Selecting the \(\tau\)
parameter carefully has a significant impact on the accuracy and efficiency
of the solution. In this work we plan to use the algorithmic version proposed
by Shakib et al. \cite{shakib-1991}.

To solve the linear system numerically, the integrals in the weak form are
approximated using Gauss-Legendre quadrature to form a stiffness matrix. By
using piecewise linear triangular finite element basis functions, the matrix
entry for each mesh node can be formed from only that node and its neighbors.
This coupling results in a sparse matrix that can be solved efficiently.

To couple nodes on part boundaries with neighbors on other processors, we
will share one ring of elements on each boundary, referred to as ``ghost''
cells. These ghost cells are supported natively by the PUMI library we plan
to use for meshing. MPI is used to communicate ghost cell data between
processors. Each processor will then use the GPU to solve the linear system.
We will demonstrate strong scaling of our parallel algorithm by obtaining a
solution for a large mesh on an increasing number of compute nodes.
Then we will demonstrate weak scaling of our algorithm by increasing the
problem size at the same rate as computate resources.

\printbibliography

\end{document}
